\chapter{Del píxel al foco: cómo se detecta un fuego}\label{cap:deteccion}

\objetivos{
\begin{itemize}
  \item Plantear el modelo del píxel mixto y calcular la temperatura de brillo de un píxel que
  contiene un fuego pequeño.
  \item Explicar por qué el tamaño mínimo detectable es proporcional al área del píxel.
  \item Describir, paso a paso, el algoritmo contextual de detección.
  \item Interpretar el valor de confianza de MODIS y de VIIRS.
\end{itemize}}

\section{El problema}

Un fuego casi nunca llena un píxel. Una quema de rastrojo ocupa unas decenas o centenas de
metros cuadrados dentro de un píxel de un millón. La pregunta de este capítulo es cómo puede un
sensor notar algo que ocupa una diezmilésima parte de lo que está mirando, y cómo decide un
programa, sin intervención humana y para todo el planeta varias veces al día, qué píxeles tienen
fuego.

\section{El píxel mixto}

El detector no distingue las partes del píxel: suma toda la radiación que llega de su huella. Si
una fracción $p$ del área está ocupada por fuego a temperatura $T_f$ y el resto, $1-p$, por un
fondo a temperatura $T_b$, la radiancia que mide en un canal de longitud de onda $\lambda$ es el
promedio ponderado por área:
\begin{equation}\label{eq:mixto}
  L_\lambda = p\,B_\lambda(T_f) + (1-p)\,B_\lambda(T_b).
\end{equation}
Este es el \termino{modelo del píxel mixto} (figura~\ref{fig:mixto}). Supone cuerpos negros y
desprecia la atmósfera, pero recoge lo esencial.

\begin{figure}[ht]
\centering
\begin{tikzpicture}
  \fill[bosque!30] (0,0) rectangle (4.5,4.5);
  \draw[ceniza,thick] (0,0) rectangle (4.5,4.5);
  \shade[inner color=yellow,outer color=fuego] (3.1,1.2) rectangle (3.55,1.65);
  \node[font=\small,align=center] at (2.25,3.75) {fondo\\ $T_b$, fracción $1-p$};
  \draw[<-] (3.33,1.7)--(3.33,2.3) node[above,font=\footnotesize,align=center]{fuego\\ $T_f$, fracción $p$};
  \draw[<->] (0,-0.3)--(4.5,-0.3) node[midway,below,font=\footnotesize]{un píxel};
  \node[font=\Large] at (5.6,2.25) {$\Longrightarrow$};
  \node[caja,align=left,text width=6.3cm,anchor=west] at (6.5,2.25) {%
    El sensor mide una sola radiancia por canal:\\[4pt]
    $L_4 = p\,B_4(T_f) + (1-p)\,B_4(T_b)$\\[2pt]
    $L_{11} = p\,B_{11}(T_f) + (1-p)\,B_{11}(T_b)$\\[4pt]
    y las convierte en dos temperaturas de brillo, $T_4$ y $T_{11}$.};
\end{tikzpicture}
\caption{El modelo del píxel mixto: una pequeña fracción muy caliente dentro de un fondo
uniforme.}
\label{fig:mixto}
\end{figure}

\begin{ejemplo}{un fuego de 100 metros cuadrados en un píxel de MODIS}
Un píxel de \SI{1}{\kilo\metre\squared} con fondo a 300~K contiene un fuego de
\SI{100}{\metre\squared} a 1000~K. Calcula las temperaturas de brillo a 3,96 y a 11~\um.
\solucion La fracción de fuego es $p = 100/10^6 = 10^{-4}$. Con las radiancias del
capítulo~\ref{cap:radiacion} ($B_4$: 0,673 para el fondo y 3320 para el fuego; $B_{11}$: 9,57 y
274):
\begin{align*}
  L_4 &= 10^{-4}\times 3320 + 0{,}9999\times 0{,}673 = 0{,}332 + 0{,}673 = 1{,}005,\\
  L_{11} &= 10^{-4}\times 274 + 0{,}9999\times 9{,}573 = 0{,}027 + 9{,}572 = 9{,}599.
\end{align*}
Convirtiendo con la ecuación~\eqref{eq:tb}: $T_4 = \SI{310.3}{\kelvin}$ y $T_{11} = \SI{300.2}{\kelvin}$.

Un fuego que ocupa una diezmilésima del píxel sube la temperatura de brillo a 4~\um{} en
\textbf{10 kelvin} y la de 11~\um{} en apenas \textbf{0,2 kelvin}. A 4~\um, la contribución del
fuego (0,332) es la mitad de la del fondo entero (0,673).
\end{ejemplo}

La figura~\ref{fig:mixto-curvas} generaliza el ejemplo a cualquier fracción de fuego.

\begin{figure}[ht]
\centering
\begin{tikzpicture}
\begin{axis}[width=13cm,height=7.2cm,xmode=log,xmin=1e-5,xmax=1e-1,ymin=295,ymax=520,
  xlabel={Fracción del píxel ocupada por fuego a 1000 K},
  ylabel={Temperatura de brillo del píxel (K)},
  legend style={at={(0.5,1.04)},anchor=south,font=\small,draw=none,legend columns=2,
    /tikz/every even column/.append style={column sep=10pt}},
  grid=major,grid style={ceniza!15},samples=120,domain=1e-5:1e-1,
  declare function={
    planck(\l,\T)=1.191e8/(\l^5*(exp(14388/(\l*\T))-1));
    tb(\l,\L)=14388/(\l*ln(1+1.191e8/(\l^5*\L)));}]
\addplot[fuego,very thick] {tb(3.96, x*planck(3.96,1000)+(1-x)*planck(3.96,300))};
\addplot[cielo,very thick] {tb(11, x*planck(11,1000)+(1-x)*planck(11,300))};
\legend{$T_4$ (3{,}96~\um),$T_{11}$ (11~\um)}
\draw[ceniza,dashed] (axis cs:1e-4,295)--(axis cs:1e-4,520);
\draw[ceniza,dashed] (axis cs:1e-3,295)--(axis cs:1e-3,520);
\node[font=\footnotesize,rotate=90,anchor=south east] at (axis cs:1e-4,518) {100 m$^2$ en 1 km$^2$};
\node[font=\footnotesize,rotate=90,anchor=south east] at (axis cs:1e-3,518) {1000 m$^2$ en 1 km$^2$};
\draw[<->,thick] (axis cs:1e-3,301.9)--(axis cs:1e-3,351.7);
\node[font=\footnotesize,anchor=west,fill=white,inner sep=1pt] at (axis cs:1.15e-3,327) {$\Delta T \approx 50$ K};
\end{axis}
\end{tikzpicture}
\caption{Temperaturas de brillo de un píxel con fondo a 300~K según la fracción ocupada por un
fuego a 1000~K, calculadas con las ecuaciones~\eqref{eq:mixto} y~\eqref{eq:tb}. El canal de
4~\um{} reacciona a fracciones minúsculas; el de 11~\um{} casi no se entera hasta que el fuego
ocupa cerca del 1\,{}\% del píxel.}
\label{fig:mixto-curvas}
\end{figure}

Dos conclusiones de la figura.

\textbf{La firma del fuego es la separación entre los dos canales.} Un píxel sin fuego tiene
$T_4$ y $T_{11}$ parecidas (de día, $T_4$ es algo mayor por la luz solar reflejada). Un píxel
con fuego tiene $T_4$ mucho mayor que $T_{11}$. Por eso la variable clave es
\begin{equation}
  \Delta T = T_4 - T_{11}.
\end{equation}

\textbf{Dos incógnitas, dos ecuaciones.} El modelo tiene dos incógnitas del fuego, $p$ y $T_f$,
y dos mediciones, $L_4$ y $L_{11}$. En principio pueden despejarse el tamaño y la temperatura
del fuego. Es el \emph{método biespectral}, propuesto por Dozier en 1981. En la práctica es muy
sensible a los errores (como el canal de 11~\um{} apenas cambia, cualquier imprecisión en el
fondo arruina la solución) y no se usa para producir los datos habituales, pero es el origen
conceptual de todo el campo.

\section{El tamaño mínimo detectable depende del píxel}

Supón que el algoritmo exige que $T_4$ supere al fondo en cierto margen, digamos 10~K. Para un
fondo a 300~K eso equivale a que la radiancia a 4~\um{} suba de 0,673 a 1,00, es decir, un
exceso de $\Delta L = 0{,}33$ unidades. Del modelo del píxel mixto, despreciando $p$ frente a 1:
\begin{equation}
  \Delta L \approx p\,B_4(T_f)
  \quad\Longrightarrow\quad
  p_{\min} \approx \frac{\Delta L}{B_4(T_f)} .
\end{equation}
La \emph{fracción} mínima no depende del tamaño del píxel. Pero el \emph{área} mínima del fuego
sí:
\begin{equation}\label{eq:amin}
  A_{\min} = p_{\min}\times A_{\text{píxel}} .
\end{equation}

\begin{definicion}{tamaño mínimo detectable}
El tamaño mínimo detectable es el área del fuego más pequeño que produce una anomalía
suficiente para que el algoritmo lo acepte. Es proporcional al área del píxel y depende
fuertemente de la temperatura del fuego y de las condiciones del fondo.
\end{definicion}

\begin{ejemplo}{el mismo fuego en MODIS y en VIIRS}
Con el criterio de 10~K sobre un fondo a 300~K, calcula el área mínima detectable de un fuego a
1000~K y de una brasa a 600~K para: (a) un píxel de MODIS en el nadir, 1~km$^2$; (b) un píxel de
MODIS en el borde, $2{,}0\times4{,}8$~km; (c) un píxel de VIIRS en el nadir, $375\times375$~m.
\solucion Las fracciones mínimas son $p_{\min} = 0{,}33/3320 = \num{1.0e-4}$ para 1000~K y
$p_{\min} = 0{,}33/287{,}5 = \num{1.1e-3}$ para 600~K (con $B_4(600) = 287{,}5$). Multiplicando
por el área de cada píxel:
\begin{center}
\begin{tabular}{lrrr}
\toprule
Píxel & Área & Fuego a 1000 K & Brasa a 600 K \\
\midrule
MODIS, nadir & \SI{1.0}{\kilo\metre\squared} & \SI{100}{\metre\squared} & \SI{1100}{\metre\squared} \\
MODIS, borde & \SI{9.6}{\kilo\metre\squared} & \SI{960}{\metre\squared} & \SI{11000}{\metre\squared} \\
VIIRS, nadir & \SI{0.14}{\kilo\metre\squared} & \SI{14}{\metre\squared} & \SI{160}{\metre\squared} \\
\bottomrule
\end{tabular}
\end{center}
Un fuego de \SI{50}{\metre\squared} a 1000~K (una quema de 7 por 7 metros) es detectable por
VIIRS y no por MODIS. Y una brasa a 600~K necesita ser diez veces más extensa que una llama a
1000~K para detectarse.
\end{ejemplo}

\begin{atencion}[Estos números son ilustrativos]
El criterio real no es un umbral fijo de 10~K sino un conjunto de pruebas, y el umbral efectivo
cambia con la variabilidad del fondo, la hora y la geometría. Los valores de la tabla sirven
para comparar sensores, no como especificación. Lo que sí es sólido es la proporcionalidad de la
ecuación~\eqref{eq:amin}.
\end{atencion}

\begin{parareto}
La ecuación~\eqref{eq:amin} es la explicación física, en una línea, de por qué VIIRS detecta más
que MODIS, y por qué la diferencia es mayor para quemas pequeñas y de baja temperatura. Es un
buen argumento para el pitch: la razón VIIRS/MODIS \emph{tiene que} depender del tipo de fuego,
porque depende de cuántos fuegos caen entre los dos límites de detección.
\end{parareto}

\section{El algoritmo contextual}

Detectar un fuego sería sencillo si el fondo fuera siempre igual: bastaría un umbral fijo. Pero
el fondo va desde un glaciar a 260~K hasta un desierto a 330~K, y de día la luz reflejada lo
altera todo. Un umbral único daría miles de falsas alarmas en el Sahara y dejaría pasar la
mayoría de los fuegos de Siberia.

La solución es comparar cada píxel con sus vecinos.

\begin{definicion}{algoritmo contextual}
Un algoritmo contextual decide si un píxel contiene fuego comparando sus temperaturas de brillo
con las de los píxeles de su entorno, que se toman como estimación de cómo se vería ese píxel
si no ardiera.
\end{definicion}

El algoritmo de MODIS procesa cada píxel en la secuencia de la figura~\ref{fig:algoritmo}. El de
VIIRS de 375~m sigue el mismo esquema con las bandas I4 e I5 y umbrales adaptados.

\begin{figure}[ht]
\centering
\begin{tikzpicture}[node distance=7mm,
  paso/.style={caja,text width=5.6cm,fill=cieloClaro},
  fin/.style={caja,fill=cenizaClaro,text width=2.9cm},
  si/.style={caja,fill=fuegoClaro,text width=2.9cm,font=\small\bfseries}]
  \node[paso] (a) {\textbf{1. Máscaras}\\ ¿es nube o agua?};
  \node[paso,below=of a] (b) {\textbf{2. Candidato}\\ ¿$T_4$ y $\Delta T$ superan un mínimo?};
  \node[paso,below=of b] (c) {\textbf{3. Prueba absoluta}\\ ¿$T_4$ es altísima?};
  \node[paso,below=of c] (d) {\textbf{4. Fondo}\\ estadísticas de los vecinos válidos};
  \node[paso,below=of d] (e) {\textbf{5. Pruebas contextuales}\\ ¿destaca sobre su fondo?};
  \node[paso,below=of e] (f) {\textbf{6. Rechazos}\\ ¿brillo solar, desierto, costa?};
  \node[si,below=of f] (g) {FUEGO\\ + confianza};
  \node[fin,right=14mm of a] (a2) {descartado};
  \node[fin,right=14mm of b] (b2) {sin fuego};
  \node[fin,right=14mm of e] (e2) {sin fuego};
  \node[fin,right=14mm of f] (f2) {sin fuego};
  \node[fin,right=14mm of d] (d2) {<<desconocido>>};
  \draw[flujo] (a)--(b) node[midway,left,font=\footnotesize]{no};
  \draw[flujo] (b)--(c) node[midway,left,font=\footnotesize]{sí};
  \draw[flujo] (c)--(d) node[midway,left,font=\footnotesize]{no};
  \draw[flujo] (d)--(e);
  \draw[flujo] (e)--(f) node[midway,left,font=\footnotesize]{sí};
  \draw[flujo] (f)--(g) node[midway,left,font=\footnotesize]{no};
  \draw[flujo] (a)--(a2) node[midway,above,font=\footnotesize]{sí};
  \draw[flujo] (b)--(b2) node[midway,above,font=\footnotesize]{no};
  \draw[flujo] (e)--(e2) node[midway,above,font=\footnotesize]{no};
  \draw[flujo] (f)--(f2) node[midway,above,font=\footnotesize]{sí};
  \draw[flujo] (d)--(d2) node[midway,above,font=\footnotesize,align=center]{pocos\\[-2pt] vecinos};
  \draw[flujo] (c.west) -- ++(-1.0,0) |- (f.west) node[pos=0.25,left,font=\footnotesize]{sí};
\end{tikzpicture}
\caption{Esquema del algoritmo contextual de detección de fuego activo.}
\label{fig:algoritmo}
\end{figure}

\subsection{Paso 1: máscaras de nubes y de agua}

Los píxeles nubosos se identifican con pruebas sencillas de reflectancia y de temperatura: una
nube es muy brillante en el visible y fría a 12~\um. Los píxeles de agua se toman de un mapa
fijo. Ninguno de los dos se evalúa como posible fuego ni se usa como vecino <<de fondo>>.

\subsection{Paso 2: candidatos}

Para no gastar cálculo en los millones de píxeles que evidentemente no arden, se descartan los
que no alcanzan unos mínimos muy permisivos. En MODIS, de día, un píxel es
\termino{candidato a fuego} si
\begin{equation}
  T_4 > \SI{310}{\kelvin}, \qquad \Delta T > \SI{10}{\kelvin}, \qquad \rho_{0{,}86} < 0{,}3,
\end{equation}
donde $\rho_{0{,}86}$ es la reflectancia en el infrarrojo cercano (un valor alto delata nube o
brillo solar). De noche, el umbral de $T_4$ baja a 305~K y no se usa la reflectancia.

\subsection{Paso 3: prueba absoluta}

Si el píxel está tan caliente que no cabe otra explicación, se acepta sin mirar el entorno. En
MODIS: $T_4 > \SI{360}{\kelvin}$ de día o $T_4 > \SI{320}{\kelvin}$ de noche.

\subsection{Paso 4: caracterización del fondo}

Para el resto de candidatos se abre una ventana cuadrada centrada en el píxel y se reúnen los
vecinos \termino{válidos}: los que no son nube, ni agua, ni a su vez posibles fuegos. La ventana
empieza pequeña y crece (hasta $21\times21$ píxeles en MODIS) hasta juntar al menos una cuarta
parte de vecinos válidos y un mínimo de ocho. Si ni así se consigue, el píxel queda clasificado
como <<desconocido>> (figura~\ref{fig:ventana}).

Con los vecinos válidos se calculan la media y una medida de dispersión de $T_4$, de $T_{11}$ y
de $\Delta T$. Se usa la \termino{desviación absoluta media}, $\delta$, que es el promedio de
las distancias de cada valor a la media; es menos sensible a valores extremos que la desviación
estándar. Llamaremos $\overline{T_4}$, $\delta_4$, $\overline{\Delta T}$ y $\delta_{\Delta T}$ a
los valores del fondo.

\begin{figure}[ht]
\centering
\begin{tikzpicture}[x=0.62cm,y=0.62cm]
  \foreach \i in {0,...,8}{\foreach \j in {0,...,8}{\fill[bosque!22] (\i,\j) rectangle ++(1,1);}}
  \foreach \i/\j in {0/8,1/8,2/8,0/7,1/7,2/7,3/8,0/6,1/6}{\fill[white] (\i,\j) rectangle ++(1,1);}
  \foreach \i/\j in {6/0,7/0,8/0,7/1,8/1,8/2}{\fill[cielo!45] (\i,\j) rectangle ++(1,1);}
  \foreach \i/\j in {5/4,5/5,3/3}{\fill[brasa!60] (\i,\j) rectangle ++(1,1);}
  \fill[fuego] (4,4) rectangle ++(1,1);
  \foreach \i in {0,...,9}{\draw[ceniza!70,thin] (\i,0)--(\i,9); \draw[ceniza!70,thin] (0,\i)--(9,\i);}
  \draw[black,very thick] (4,4) rectangle ++(1,1);
  \draw[black,thick,dashed] (2,2) rectangle (7,7);
  \begin{scope}[shift={(10.3,7.6)}]
    \fill[fuego] (0,0) rectangle ++(0.7,0.7); \node[right,font=\footnotesize] at (0.9,0.35) {píxel que se evalúa};
    \fill[bosque!22] (0,-1.1) rectangle ++(0.7,0.7); \node[right,font=\footnotesize] at (0.9,-0.75) {vecino válido (fondo)};
    \fill[brasa!60] (0,-2.2) rectangle ++(0.7,0.7); \node[right,font=\footnotesize] at (0.9,-1.85) {otro posible fuego: excluido};
    \fill[white] (0,-3.3) rectangle ++(0.7,0.7); \draw[ceniza] (0,-3.3) rectangle ++(0.7,0.7); \node[right,font=\footnotesize] at (0.9,-2.95) {nube: excluida};
    \fill[cielo!45] (0,-4.4) rectangle ++(0.7,0.7); \node[right,font=\footnotesize] at (0.9,-4.05) {agua: excluida};
    \draw[black,thick,dashed] (0,-5.5) rectangle ++(0.7,0.7); \node[right,font=\footnotesize] at (0.9,-5.15) {ventana actual ($5\times5$)};
  \end{scope}
\end{tikzpicture}
\caption{Ventana de fondo. Solo los vecinos válidos entran en las estadísticas. Si no hay
suficientes, la ventana se agranda.}
\label{fig:ventana}
\end{figure}

\subsection{Paso 5: pruebas contextuales}

El candidato se acepta como fuego si destaca de su fondo en las dos variables. Las condiciones
centrales del algoritmo de MODIS son
\begin{align}
  \Delta T &> \overline{\Delta T} + 3{,}5\,\delta_{\Delta T}, \label{eq:ctx1}\\
  \Delta T &> \overline{\Delta T} + \SI{6}{\kelvin}, \label{eq:ctx2}\\
  T_4 &> \overline{T_4} + 3\,\delta_4 . \label{eq:ctx3}
\end{align}
La primera y la tercera son relativas: exigen que el píxel se aparte de la media varias veces lo
que suelen apartarse sus vecinos. En un entorno homogéneo ($\delta$ pequeña) basta una anomalía
modesta; en uno heterogéneo se exige más. La segunda impone un mínimo absoluto de 6~K para que,
en fondos extremadamente uniformes, una diferencia insignificante no dispare una detección. De
día hay condiciones adicionales sobre el canal de 11~\um.

\begin{ejemplo}{aplicar las pruebas contextuales}
Un píxel diurno tiene $T_4 = \SI{318.0}{\kelvin}$ y $T_{11} = \SI{299.5}{\kelvin}$. Sus vecinos
válidos tienen $\overline{T_4} = \SI{304.0}{\kelvin}$, $\delta_4 = \SI{1.8}{\kelvin}$,
$\overline{\Delta T} = \SI{5.5}{\kelvin}$ y $\delta_{\Delta T} = \SI{1.2}{\kelvin}$. ¿Es fuego?
\solucion $\Delta T = 318{,}0 - 299{,}5 = \SI{18.5}{\kelvin}$. Es candidato: $T_4 > 310$ y
$\Delta T > 10$. No pasa la prueba absoluta ($318 < 360$). Pruebas contextuales:
\begin{itemize}
  \item \eqref{eq:ctx1}: $5{,}5 + 3{,}5\times1{,}2 = 9{,}7$; $18{,}5 > 9{,}7$. Se cumple.
  \item \eqref{eq:ctx2}: $5{,}5 + 6 = 11{,}5$; $18{,}5 > 11{,}5$. Se cumple.
  \item \eqref{eq:ctx3}: $304{,}0 + 3\times1{,}8 = 309{,}4$; $318{,}0 > 309{,}4$. Se cumple.
\end{itemize}
El píxel se clasifica como fuego, salvo que lo rechace alguna prueba del paso 6.

Si el mismo píxel estuviera en un paisaje muy heterogéneo con $\delta_{\Delta T} =
\SI{4.5}{\kelvin}$, la primera condición exigiría $\Delta T > 5{,}5 + 15{,}75 = 21{,}3$~K y el
píxel \emph{no} sería fuego. El mismo fuego se detecta o no según cómo sea su entorno.
\end{ejemplo}

\subsection{Paso 6: rechazo de falsas alarmas}

Por último se descartan las detecciones diurnas que coinciden con las situaciones del
capítulo~\ref{cap:atmosfera}: geometría de brillo solar, límites de desierto (cuando muchos
vecinos quedaron excluidos por ser candidatos sin serlo realmente) y bordes de costa.

\section{Confianza}

Cada detección lleva una estimación de cuán segura es.

En \textbf{MODIS} es un número entre 0 y 100. Se calcula combinando cuánto supera el píxel los
umbrales: qué tan alto es $T_4$, cuántas desviaciones se aparta del fondo en $T_4$ y en
$\Delta T$, y cuántos vecinos son nube o agua. Por convención se agrupa en tres clases:
\emph{baja} (menos de 30), \emph{nominal} (de 30 a 80) y \emph{alta} (más de 80).

En \textbf{VIIRS} de 375~m es directamente una de tres categorías:
\begin{itemize}
  \item \emph{low} (\texttt{l}): detecciones diurnas en zona de posible brillo solar o con una
  anomalía térmica débil;
  \item \emph{nominal} (\texttt{n}): detecciones sin esas circunstancias y con una anomalía
  clara;
  \item \emph{high} (\texttt{h}): píxeles saturados, de día o de noche.
\end{itemize}

\begin{atencion}[La confianza no es una probabilidad y no es comparable entre sensores]
Un valor de 70 en MODIS no significa <<70\,{}\% de probabilidad de ser fuego>>. Y la clase <<alta>>
de VIIRS (píxel saturado) no equivale a la clase <<alta>> de MODIS (número mayor que 80). Si
filtras por confianza, define el filtro por separado para cada sensor y mantén ese criterio fijo
en todo el periodo. Un criterio habitual y fácil de defender: descartar solo la clase baja en
ambos.
\end{atencion}

\section{Diferencias del algoritmo de VIIRS}

El algoritmo de 375~m hereda el esquema de MODIS, con adaptaciones a su instrumento:

\begin{itemize}
  \item usa I4 (3,74~\um) e I5 (11,45~\um), más las bandas I1 a I3 para nubes, agua y brillo
  solar;
  \item como I4 se satura a 367~K, trata de forma especial los píxeles saturados, e incluso los
  casos en que la saturación extrema produce valores anómalamente bajos;
  \item ajusta sus umbrales de candidatura según las condiciones de la escena;
  \item presta atención a la región de la Anomalía del Atlántico Sur, donde la radiación del
  espacio produce ruido en el detector que puede simular fuegos nocturnos. Esa región incluye
  buena parte de Sudamérica.
\end{itemize}

El resultado es un producto con errores de comisión bajos (del orden del 1\,{}\% para la clase
nominal en la evaluación inicial) y una sensibilidad a fuegos pequeños muy superior a la de
MODIS, sobre todo de noche.

\resumen

\begin{itemize}
  \item El \textbf{modelo del píxel mixto} combina fuego y fondo ponderados por área:
  $L = p\,B(T_f) + (1-p)\,B(T_b)$.
  \item Un fuego de 100~m$^2$ a 1000~K en un píxel de 1~km$^2$ eleva $T_4$ unos 10~K y $T_{11}$
  unos 0,2~K. La firma del fuego es $\Delta T = T_4 - T_{11}$.
  \item El \textbf{área mínima detectable} es proporcional al área del píxel: unas siete veces
  menor en VIIRS que en MODIS en el nadir, y la diferencia crece hacia los bordes.
  \item El \textbf{algoritmo contextual} filtra nubes y agua, selecciona candidatos, caracteriza
  el fondo con los vecinos válidos y acepta el píxel si destaca varias desviaciones sobre él.
  \item Que un fuego se detecte depende también de la homogeneidad de su entorno.
  \item La \textbf{confianza} es un número de 0 a 100 en MODIS y una categoría en VIIRS. No son
  equivalentes.
\end{itemize}

\ejercicios

\begin{enumerate}
  \item Repite el primer ejemplo del capítulo para un fuego de \SI{1000}{\metre\squared} a
  800~K ($B_4(800) = 1317$, $B_{11}(800) = 179$). Obtén $T_4$, $T_{11}$ y $\Delta T$.
  \item ¿Qué fracción de un píxel debe ocupar un fuego a 1000~K para \emph{saturar} la banda I4
  de VIIRS (367~K) con fondo a 300~K? Exprésala como área dentro de un píxel de
  $375\times375$~m. Pista: calcula primero la radiancia de un cuerpo negro a 367~K en
  3,74~\um.
  \item Usando la ecuación~\eqref{eq:amin}, ¿cuánto mayor es el fuego mínimo detectable por
  MODIS a 900~km del nadir (píxel de unos $1{,}7\times3{,}1$~km) que en el nadir?
  \item Un píxel nocturno tiene $T_4 = \SI{309}{\kelvin}$, $T_{11} = \SI{287}{\kelvin}$, y su
  fondo $\overline{T_4} = \SI{289}{\kelvin}$, $\delta_4 = \SI{0.9}{\kelvin}$,
  $\overline{\Delta T} = \SI{1.5}{\kelvin}$, $\delta_{\Delta T} = \SI{0.6}{\kelvin}$. Aplica los
  pasos 2, 3 y 5.
  \item Explica por qué se excluyen del fondo los vecinos que son a su vez candidatos a fuego.
  ¿Qué le pasaría a la detección en el centro de un incendio grande si no se excluyeran?
  \item Dos quemas idénticas ocurren a la vez: una en una pampa uniforme y otra en un mosaico de
  chacras, eriazos y bosque. ¿Cuál tiene más probabilidad de ser detectada? ¿Qué implica esto
  para comparar el número de focos entre dos paisajes distintos?
  \item ¿Por qué no es correcto filtrar <<confianza $\geq 50$>> en un archivo que mezcla
  detecciones de MODIS y de VIIRS?
\end{enumerate}

\lecturas

Giglio, Descloitres, Justice y Kaufman (2003) describen el algoritmo contextual original de
MODIS, y Giglio, Schroeder y Justice (2016) su versión actual (Colección~6). Schroeder y
colaboradores (2014) presentan el algoritmo de VIIRS de 375~m. Véase el
apéndice~\ref{ap:fuentes}.
